\documentclass[11pt]{article}
% Shared preamble for the social optimum derivations. Compile from this folder:
%   pdflatex pigou.tex   (the network diagrams are read from ../diagrams)
\usepackage[T1]{fontenc}
\usepackage{lmodern}
\usepackage{microtype}
\usepackage[margin=2.5cm]{geometry}
\usepackage{amsmath,amssymb}
\usepackage{booktabs}
\usepackage{float}
\usepackage{parskip}
\usepackage[font=small,labelfont=bf]{caption}
\usepackage[hyphens]{url}
\usepackage{standalone}
% Shared TikZ styles for the network diagrams.
\usepackage{tikz}
\usetikzlibrary{arrows.meta,positioning}
\definecolor{nodeblue}{HTML}{2A78D6}
\definecolor{endfill}{HTML}{E3EEFB}
\definecolor{shortcut}{HTML}{D64545}
\tikzset{
  v/.style={circle, draw=nodeblue, line width=1pt, minimum size=9mm, inner sep=0pt, font=\small},
  od/.style={v, fill=endfill},
  e/.style={-{Latex[length=2.6mm]}, line width=0.8pt},
  free/.style={e, shortcut},
  c/.style={font=\footnotesize, fill=white, inner sep=1.5pt},
}


\newcommand{\SCopt}{SC_{\mathrm{opt}}}
% Repository file or folder, e.g. \file{results/grid/summary.csv}; breaks at slashes.
\DeclareRobustCommand{\file}[1]{\path{#1}}


\begin{document}

\section{Social Optimum: Definition}

\paragraph{Social cost.} An edge $e$ crossed by $f$ agents in the same timestep costs each of them
$c_e(f) = a_e + b_e f$. The social cost of a run is the total paid by all agents,
\[
  SC = \sum_t \sum_e f_e(t)\, c_e\bigl(f_e(t)\bigr),
\]
where $f_e(t)$ is the number of agents crossing $e$ in step $t$.

\paragraph{Social optimum.} $\SCopt$ is the lowest social cost any routing of the $N$ agents can achieve
under the simulator's movement rules:
\begin{enumerate}
  \item all agents start at $t = 0$;
  \item every agent that has not arrived crosses exactly one edge per step (no waiting);
  \item no agent moves back towards its origin or revisits a node (\texttt{forward\_only}).
\end{enumerate}
Agents are indivisible. Edge costs do not change with $N$, so $\SCopt$ is computed separately for each
population size and reported per agent, $\SCopt/N$.

\paragraph{As a choice of paths.} Under these rules an agent crosses the $k$-th edge of its path at step
$k - 1$. A routing is therefore just a path for each agent, and
\[
  \SCopt = \min_{\text{paths}} \sum_{e,\,t} X_{e,t}\,\bigl(a_e + b_e X_{e,t}\bigr),
\]
where $X_{e,t}$ is the number of agents whose $(t{+}1)$-th edge is $e$. Agents congest each other only if
they cross the same edge in the same step.

\paragraph{Relation to the static optimum.} The classical static optimum is
$\min \sum_e x_e\, c_e(x_e)$ over the same paths, where $x_e$ is the total number of agents using $e$:
every user of an edge is charged for all of its users, whenever they cross.
\begin{itemize}
  \item If every edge is always crossed at the same step (Pigou, Diamond, Grid), then $X_{e,t} = x_e$ and
        the two optima are equal.
  \item Otherwise $\SCopt$ is never higher, since $\sum_t X_{e,t}^2 \le \bigl(\sum_t X_{e,t}\bigr)^2$:
        spreading an edge's users over several steps can only reduce congestion. For Braess and Sioux
        Falls it is strictly lower.
\end{itemize}
Every simulated run costs at least $\SCopt$, so $SC/\SCopt \ge 1$ measures how far a run is from optimal.

\end{document}
